\documentclass[11pt,a4paper]{article}
\usepackage[T1]{fontenc}
\usepackage{lmodern}
\usepackage{isabelle,isabellesym}
\usepackage{amsmath,amssymb,amsthm}
\usepackage{stmaryrd}
\usepackage{wasysym}
\usepackage{textcomp}
\usepackage[margin=1in]{geometry}
\usepackage{pdfsetup}

\isabellestyle{it}
\isadroptag{theory}
\renewcommand{\isastyletext}{\normalsize\normalfont}
\renewcommand{\isastyletxt}{\normalfont}

% bra-ket macros used in the commentary
\newcommand{\ket}[1]{|#1\rangle}
\newcommand{\bra}[1]{\langle #1|}
\newcommand{\braket}[2]{\langle #1|#2\rangle}
\newcommand{\ketbra}[2]{|#1\rangle\langle #2|}
\newcommand{\norm}[1]{\|#1\|}
\DeclareMathOperator{\Tr}{Tr}
\DeclareMathOperator{\tr}{tr}
\DeclareMathOperator{\rank}{rank}
\DeclareMathOperator{\range}{range}
\DeclareMathOperator{\spn}{span}
\DeclareMathOperator{\diag}{diag}
\DeclareMathOperator{\crk}{cr}
\newcommand{\C}{\mathbb{C}}
\newcommand{\R}{\mathbb{R}}
\newcommand{\Z}{\mathbb{Z}}
\newcommand{\Hil}{\mathcal{H}}
\newcommand{\Kil}{\mathcal{K}}

\theoremstyle{definition}
\newtheorem*{notestheorem}{Theorem (notes)}
\newtheorem*{notesdef}{Definition (notes)}
\newtheorem*{notesprop}{Proposition (notes)}
\newtheorem*{notescor}{Corollary (notes)}
\newtheorem*{noteslemma}{Lemma (notes)}
\newtheorem*{notespost}{Postulate (notes)}
\newtheorem*{notesremark}{Remark}

% math in section titles: plain-text forms for PDF bookmarks
\pdfstringdefDisableCommands{\def\rangle{>}\def\langle{<}\def\otimes{x}\def\cdots{...}%
  \def\le{<=}\def\ge{>=}\def\to{->}\def\mapsto{|->}\def\oplus{+}\def\kappa{kappa}\def\pi{pi}}
\begin{document}

\title{Lectures on Quantum Computing and Completely Positive Maps\\[0.5em]
  \large A formalisation in Isabelle/HOL of Vern Paulsen's Houston 2011 lecture notes}
\author{Dustin Bryant\\[0.3em]\small formal statements and proofs: the theories listed in the table of contents}
\maketitle

\begin{abstract}
This document is generated by Isabelle from the project theories.  Every
definition and theorem shown here has been checked by Isabelle/HOL (proofs are
folded).  The commentary restates Paulsen's lecture notes lecture by lecture
(Day~1 = 22~August~2011, Day~34 = 18~November~2011): definitions and results in
the notes' own language, each followed by the formal statement.  The informal
proofs of the notes are omitted; the formal proofs replace them.  Where the
formalisation deviates from the notes, the commentary says so.  The companion
file \texttt{NOTES\_AUDIT.md} lists, for every lecture, what is formalised and
what is not.  General matrix facts that the lectures use but do not teach form the parent
session \emph{Essential Functional Analysis}, with its own document.  Two companion entries build on it: the functional calculus of Hermitian
matrices with product formulas, and quantum phase estimation with the HHL algorithm.
\end{abstract}

\tableofcontents
\newpage

\parindent 0pt\parskip 0.5ex

\input{QC_Notes_Overview}
\part{Basic mathematics (Days 1--7)}
\input{Day01_Inner_Product}
\input{Day02_Trace_Projections}
\input{Day03_Unitary_Matrices}
\input{Day04_Hermitian_Matrices}
\input{Day05_Positivity_Tests}
\input{Day05_Sylvester_Criterion}
\input{Day06_Tensor_Products}
\input{Day07_Postulates}
\part{An introduction to quantum computing (Days 8--17)}
\input{Day08_Measurements}
\input{Day09_POVMs_Entanglement}
\input{Day09_Bell_Measurement}
\input{Day10_Quantum_Gates}
\input{Day11_Gates_No_Cloning}
\input{Day12_Quantum_Parallelism}
\input{Day13_Density_Matrices}
\input{Day14_Pure_States}
\input{Day15_Partial_Traces}
\input{Day16_Block_Form}
\input{Day17_Quantum_Noise}
\part{Completely positive maps (Days 18--23)}
\input{Day18_CP_Maps}
\input{Day19_Choi_Theorem}
\input{Day20_Positive_Functionals}
\input{Day21_Unitary_Freedom}
\input{Day21_Average_Conjecture}
\input{Day22_Minimal_Kraus}
\input{Day22_UCP_Compact}
\input{Day23_Choi_Extreme_Points}
\input{Day23_Extremality_Examples}
\part{Operator systems, Arveson's theorems and entanglement (Days 24--28)}
\input{Day24_Operator_Systems}
\input{Day25_Hahn_Banach}
\input{Day25_Arveson_Correspondence}
\input{Day25_Unital_Functional_Norm}
\input{Day26_Arveson_Extension}
\input{Day26_Separability}
\input{Day27_Detecting_Entanglement}
\input{Day28_Entanglement_Witnesses}
\part{Quantum error correction (Days 29--34)}
\input{Day29_Error_Correction}
\input{Day30_Bit_Flip_Code}
\input{Day31_Bit_Flip_Channel}
\input{Day32_Phase_Flip_Code}
\input{Day32_Recovery_Proposition}
\input{Day33_Shor_Code}
\input{Day34_Knill_Laflamme}
\part{Beyond the notes}
\input{Quantum_Channels}

\end{document}
